% Pursuing Stacks (cotes 134-2 to 134-8): archive cote 134-2, site page 64 (PDF page 65).
% appendix (letters to Breen). Draft compared: ps1.tex:1474-1497.
% AI-transcribed by Claude Opus 5.5 (claude-opus-5-5, adaptive thinking, effort high), Message Batches API, batch ps1 (msgbatch_01Rv5hCaNTtUt4JHx1yqoHPL), 2026-11-12; not yet checked by eye.
\scanpage{134-2/64}
\mspage{- \unsure{\hand{4}}9 -}

formula should result without difficulty from the usual \hand{``}dévissages\hand{''} which
allow one to verify the duality in the particular standard cases $F = \mathbb{Z}/p\mathbb{Z}$ ,
$\mu_p$ , $\alpha_p$ on a smooth, complete $X$ . (N.B. the case of coefficients prime
to the characteristic is already known.) Let us make explicit what the
formula of duality says for $R^1 f_*(\mathbb{G}_X)$ , where $G$ is a finite group étale
on $k$ (the most important case being $G = (\mathbb{Z}/p^m\mathbb{Z})_k$ ); one recovers \deleted{the} \inserted{Serre's}
description \deleted{of Serre} of ``geometric classfield\deleted{s''\unsure{,}} \inserted{theory''} in \deleted{the form of} \inserted{terms of} extensions
by $G$ of a generalised jacobian of $X$ . Thus, the duality formula can be
understood as a cohomological version, considerably enriched, of geometric
class field theory. When the base field $k$ is finite, to retrieve the
classfield theory in the classical form, one can use ``the trick of Lang''
[on the relation between the ``arithmetic'' $\pi_1$ of a smooth, connected
commutative algebraic group $J$ on $k$ and its $H^0(k, J) = J(k)$ : the $\pi_1^{ar}(J)$
classifies the isogenies above $J$ with kernel a constant group
$\pi_1^{ar}(J) \simeq H^0(k, J)$] - in its cohomological form, which may be stated:
\[ \Delta_0 \mathbb{R}\Gamma_k(J^*) \simeq \mathbb{R}\Gamma_k(\Delta J^*[1]) \hand{,} \]
where $\Delta$ is Serre duality, $\Delta_0$ Pontrjagin duality for the totally disconnected
topological abelian groups (duality with values in $\mathbb{Q}/\mathbb{Z}$)\hand{)}, $J^*$ a complex of
algebraic ind-progroups on $k$ . Taking account of th\over{e}{is} \deleted{``formula of Lang duality''} \inserted{``Lang duality formula''}
and applying $\mathbb{R}\Gamma_k$ to the formula of duality for geometric classfields, one \inserted{gets}
\deleted{finds} \inserted{the} ``duality formula \over{for}{of} arithmetic classfield\deleted{s''} \inserted{theory''}:
\[ \Delta_0(H_!(X, F)) \simeq H^*(X, D(F)[3]) \]
(isomorphism of totally disconnected topological groups).

Another remark: when $F$ \deleted{comes from an} \inserted{is not an} ``étale sheaf'', but has a
continuous structure such as $\alpha_p$ , one must be careful in the definition of
$Rf_!(F)$ , for $X$ non complete, starting from the compactification $\tilde{X}$ ; thus,
if $F$ comes from an ``admissible'' sheaf $\tilde{F}$ on $\tilde{X}$ , one must have an exact
triangle

\margin{64}
